% ===== 原文 pp.59–62（PDF p.69–72） =====
% [原文 p.59]
\chapter{4d TFT 的面算子}

在写下 $\cN=4$ SYM 的各个扭转版本并分析了它们的 3d 约化之后，我们现在可以来分类这些理论中面算子的谱。正如第 4 章所强调的，给定 4d TFT 的面算子具有 2-范畴的结构；也就是说，它是约化后的 3d TFT 的边界条件所构成的 2-范畴。

本章系统地利用这一观察来分类 $t=\ii$、$t=1$ 与 $t=0$ 时 GL-扭转理论的面算子。（根据上一章所述，Vafa-Witten 理论的面算子谱将与 $t=0$ 理论的相同。）3d 视角的好处在于：边界条件的分类问题是更为人熟悉的问题。

本章的材料在 [4] 中给出。对约化理论边界条件的大部分分析依赖于关于 2d 与 3d TFT 的工作，这些工作被放在附录 A 与 B 中。

\section{Gukov-Witten 面算子}

GL-扭转理论中最简单的面算子由 S. Gukov 与 E. Witten [2] 引入。它们是对应于规范场 $A$ 与玻色 1-形式 $\phi$ 中余维 2 奇点的无序算子，其形式为
\[
A=\alpha\dd\theta,\qquad \phi=\beta\frac{\dd r}{r}-\gamma\dd\theta .
\]
这里 $\alpha$ 是 $G$ 的极大环面 $T$ 的元素，而 $\beta,\gamma$ 是 $T$ 的李代数 $\ft$ 的元素。为简单起见，设三元组 $(\alpha,\beta,\gamma)$ 把 $G$ 破缺到 $T$。保持 $T$ 的规范变换构成 Weyl 群 $\cW$；三元组 $(\alpha,\beta,\gamma)$ 的取值要模去 $\cW$ 在 $T\times\ft\times\ft$ 上的作用。除 $A$ 与 $\phi$ 之外，所有场都非奇异。

该面算子还依赖于一个额外参数 $\eta$，它取值于环面 $\Hom(\Lambda_{\mathrm{cochar}},U(1))$。这里 $\Lambda_{\mathrm{cochar}}$ 是磁荷格 $\Hom(U(1),T)$。等价地，如 [2] 所解释，$\eta$ 可以看作取值于 $\LT$，即朗兰兹对偶群的极大环面。

% [Tutor Note: p.59 参数 $\eta$ 的取值域核对：$\Hom(\Lambda_{\mathrm{cochar}},U(1))$ 就是朗兰兹对偶群的极大环面。取 $G=U(1)$：$\Lambda_{\mathrm{cochar}}=\Hom(U(1),U(1))=\Z$，故 $\Hom(\Z,U(1))=U(1)=\LT$，可直接验算；一般情形下"磁荷格的对偶 = 对偶群的极大环面"是 $\LG$ 的标准构造（$\Lambda_{\mathrm{cochar}}$ 是 $\LG$ 的权格）。]

% [原文 p.60]
参数 $\eta$ 的来源如下。首先注意，上述场中的奇点把规范群破缺到 $T$。因此若 $D$ 是面算子所支撑的余维 2 子流形，则规范场在 $D$ 上的限制有第一陈类 $c_1|_D$，它取值于 $\Lambda_{\mathrm{cochar}}$。给定 $\eta$，我们可以往路径积分中插入一个相位因子
\[
\eta(c_1(D)).
\]
这个因子只依赖于规范场在 $D$ 上的行为，可以看作对上文所定义的面算子的一个依赖于 $\eta$ 的修正。

Gukov-Witten 面算子对任意 $t$ 都是 BRST 不变的，但它们的性质依赖于 $t$。我们将在下文看到，还存在许多别的面算子。在下文中我们将集中于 $t=\ii$、$t=1$ 与 $t=0$ 三种情形。前两种情形（在 $\theta$ 角为零时）由 S-对偶互换，并在几何朗兰兹纲领的物理进路中起突出作用 [1]。最后一种情形是自对偶的，是理解量子几何朗兰兹对偶性 [35] 最自然的出发点。

我们将发现，对 $t$ 的每个取值，都存在远比这更一般的面算子。

% [Tutor Note: p.60 相位因子 $\eta(c_1(D))$ 是良定义的：$c_1|_D$ 取值于 $\Lambda_{\mathrm{cochar}}$，而 $\eta\in\Hom(\Lambda_{\mathrm{cochar}},U(1))$，二者配对给 $U(1)$ 元素。核对时可把它看作 $\LG$ 的极大环面元素作用于 $\LG$ 的余权格，这正是把 $\eta$ 解释为朗兰兹对偶侧数据的含义（[2]）。]

% [Tutor Note: p.60 "$t=\ii$ 与 $t=1$ 由 S-对偶互换、$t=0$ 自对偶"这一断言，可用 §2.4–§2.5 的 S-对偶作用（$\tau\to-1/\tau$、$G\leftrightarrow\LG$）与扭转参数 $t$ 的变换规则来复核；注意 Vafa–Witten 理论并不等同于 $t=0$ 理论，原文只说二者面算子谱相同（p.59 开头）。]

\section{\texorpdfstring{$t=\ii$}{t=i} 且 \texorpdfstring{$G=U(1)$}{G=U(1)} 时的面算子}

在 §5.5 中我们发现，$W\times S^1$ 上的 $t=\ii$ 理论约化为 $W$ 上以 $T^*G_\C$ 为靶空间的带规范场 Rozansky-Witten 模型。我们先取 $G=U(1)$，分析以 $T^*\C^*$ 为靶空间的带规范场 RW 模型的边界条件。边界条件的分析可以干净地分成物质扇区的边界条件与规范扇区的边界条件；我们依次讨论。

% [Tutor Note: p.60 此处的约化是 §5.5 的结果（圆紧化 + 3d 极限）。$G=U(1)$ 时 $T^*G_\C=T^*\C^*$ 的复维数为 2，所以其复拉格朗日子流形复维数必须是 1；§6.2.1 中"只有零截面与纤维两类"的结论即由此定维而来。]

% [Tutor Note: p.60 为何可以分两个扇区处理：带规范场 RW 模型 = RW 物质扇区 + B-型 3d 规范扇区，二者只通过边界上 $\tau$ 的边界值（即 $h=\exp(-2\pi\tau)$）与 $\cA$ 的边界值耦合（见 p.63 末段：纤维化在 $\C^*$ 上决定与 $\tau$ 的耦合，$\C^*$ 作用决定与 $\cA$ 的耦合）。因此可以先分别分类（§6.2.1、§6.2.2），再于 §6.2.3 把两面拼起来。]

\subsection{RW 模型的边界条件，靶空间 \texorpdfstring{$T^*\C^*$}{T*C*}}

根据 [36]，RW 模型中最简单的边界条件对应于 $X$ 的复拉格朗日子流形。若要保住鬼数对称性，这些拉格朗日子流形必须对伸缩 $\sigma\mapsto\lambda^2\sigma$（$\lambda\in\C^*$）不变。这要求 $T^*\C^*$ 的拉格朗日子流形是 $\C^*$ 中某个复子流形的余法丛。这意味着一个（闭的）$\C^*$-不变的复拉格朗日子流形只能是零截面 $\sigma=0$，或者是余切丛的纤维 $\tau=\tau_0$ 之一。零截面边界条件有特殊地位，在本小节中记作 $\mathbb{X}_0$。

% [Tutor Note: 上述分类可这样复核：$T^*\C^*$ 复维为 2，复拉格朗日子流形复维须为 1。$\C^*$ 的复子流形只有两种：整个 $\C^*$（其余法丛即零截面 $\sigma=0$，复维 1）与单点（其余法丛即纤维 $\tau=\tau_0$，复维 1）。既要求是复子流形的余法丛，就排除其他复维 1 的子流形（如 $\C^*$ 中的全纯曲线、$\sigma$ 与 $\tau$ 混合定义的截面）。注意本段中的 $X$ 指 RW 模型的靶空间，在本小节即 $X=T^*\C^*$。]

% [原文 p.61]
更一般的边界条件对应于由某个复拉格朗日子流形上的点所参数化的一族 B-模型或 Landau-Ginzburg 模型。如 [36] 所述、并在 [37] 中更详细地解释的，只需把自己限制在拉格朗日子流形为零截面 $\sigma=0$ 的情形即可。可以更代数地描述这些边界条件如下。回忆：边界 $\mathbb{X}_0$ 上的边界线算子范畴是一个幺半范畴，我们记作 $\mathbb{V}_{\mathbb{X}_0\mathbb{X}_0}$。给定任意边界条件 $\mathbb{X}$，可以考虑把 $\mathbb{X}$ 与 $\mathbb{X}_0$ 分开的边界缺陷线所构成的范畴 $\mathbb{V}_{\mathbb{X}\mathbb{X}_0}$。这个范畴是幺半范畴 $\mathbb{V}_{\mathbb{X}_0\mathbb{X}_0}$ 上的模范畴。[36] 提出，这个模范畴完全刻画了边界条件 $\mathbb{X}$。具体地，在以 $T^*[2]\C^*$ 为靶空间的 RW 模型中，边界线算子范畴 $\mathbb{V}_{\mathbb{X}_0\mathbb{X}_0}$ 等价于 $D^b(Coh(\C^*))$。看到这一点的一种方式是：把 RW 模型在以边界条件 $\mathbb{X}_0$ 为两端的区间上约化。得到的 2d TFT 是以 $\C^*$ 为靶空间的 B-模型，其膜的范畴可以等同于 $D^b(Coh(\C^*))$。2d 视角无法确定幺半结构，但可以证明它由通常的导出张量积给出 [36,37]。

进一步地，[36,37] 论证了：以 $T^*[2]\C^*$ 为靶空间的 RW 模型的边界条件 2-范畴，等价于 $D^b(Coh(\C^*))$ 上的模范畴所构成的导出 2-范畴。也就是说，它就是 B. Toen 与 G. Vezzosi [44] 所定义的 $\C^*$ 上导出范畴层的 2-范畴。这给出了所有边界条件的边界线算子及其 OPE 的代数描述。

% [Tutor Note: p.61 记号 $T^*[2]\C^*$ 中的 "[2]" 指余切丛的纤维坐标（即 $\sigma$）被形式地赋以鬼数 2；这一点在 p.65–p.66 被反复使用（例如 $\mathbb{C}[2]$、$N^*[2]\mathcal{C}$）。因此"伸缩 $\sigma\mapsto\lambda^2\sigma$"是 $\C^*$ 的权 2 作用，与鬼数相容。]

\subsection{\texorpdfstring{B-型 3d 规范理论的边界条件（$G=U(1)$）}{B-型 3d 规范理论的边界条件（G=U(1)）}}

回忆 §5.5：玻色场 $A|_W$、$\phi|_W$ 与费米子场 $\psi+\ii\tilde\psi,\chi,\eta-\ii\tilde\eta,\psi_4-\ii\tilde\psi_4$ 正确地约化为 B-型 3d 规范理论（在 §B.4 中讨论）的场，其作用量由下式给出：
\[
S=-\frac{1}{2e^2}\delta_Q\int_W\left(\chi\wedge\star\cF-\frac{\ii}{2}(\eta-\ii\tilde\eta)\wedge\star\dd^{\star}\phi\right)
+\frac{1}{2e^2}\int_W(\psi_4-\ii\tilde\psi_4)\dd\chi .
\]
原则上我们应当对理论做规范固定并相应地修改 BRST 算子，但在整篇论文中我们略去这一复杂性。

% [Tutor Note: p.61 作用量中 $d^{\star}$ 是余微分 $d^{\star}=\star d\star$（附录 B 中明确采用此约定），故 $\star d^{\star}\phi$ 是 3-形式，才能与 0-形式的 $(\eta-\ii\tilde\eta)$ 作楔积并在 $W$ 上积分；若误读成 $\star d\star\phi$（0-形式）则量纲不自洽。核对时先数形式次数：$\chi$ 为 2-形式、$\star\cF$ 为 1-形式，第一项为 3-形式。]

在任何规范理论中，最自然的边界条件都是 Dirichlet 与 Neumann 条件。Dirichlet 条件要求 $A+\ii\phi$ 在边界上的限制为平凡。此外还要求 $\phi_3$（即 $\phi$ 垂直于边界的分量）满足 Neumann 条件 $\partial_3\phi_3=0$。BRST 不变性随后固定了费米子的边界条件：形式 $\psi+\ii\tilde\psi$、$\chi$ 与 $\eta-\ii\tilde\eta$ 在边界上的限制必须为零。

% [Tutor Note: p.61 原文此处作"must vanish, The Neumann"，逗号后接大写 The，应为句号断句（"must vanish. The Neumann"）。译文按句意断为两句；这是原文的排印笔误，不影响内容。]

Neumann 边界条件不约束 $\cA$ 在边界上的限制，但要求 1-形式 $\star\cF=\star\dd\cA$ 的限制为零。此外 $\phi_3$ 必须满足 Dirichlet 边界条件，即它必须在边界上取一个给定值。在 Neumann 情形，BRST 不变性要求费米子 $\star\chi$、$\psi_3+\ii\tilde\psi_3$ 与 $\psi_4-\ii\tilde\psi_4$ 的限制为零。注意在 Dirichlet 情形规范群
% [原文 p.62]
在边界上被完全破缺，而在 Neumann 情形则未被破缺。

Dirichlet 条件没有任何参数，而 Neumann 条件似乎依赖于一个实参数 $\beta$，即 $\phi_3$ 的边界值。在量子层面上还有另一个参数：我们可以往作用量中加一个边界拓扑项
\[
\theta\int_{\partial W}\frac{\cF}{2\pi}.
\]
事实上这两个参数都是无关的，意思是拓扑关联函数不依赖于它们。参数 $\theta$ 的无关性来自一个事实：上述拓扑项是 BRST-正合的，且等于
\[
\frac{\theta}{2\pi}\delta_Q\int_{\partial W}\chi .
\]
要看出参数 $\beta$ 的无关性，注意：平移 $\beta$ 需要往作用量中加一个正比于
\[
\int_{\partial W}\partial_3\phi_3
\]
的边界项。由于 $\phi_1$ 与 $\phi_2$ 在边界上为零，这一项也是 BRST-正合的，且正比于
\[
\delta_Q\int_{\partial W}(\eta-\ii\tilde\eta).
\]

% [Tutor Note: p.62 第一个无关性论证的核对方法：原文说拓扑项是 BRST-正合的，且等于 $\frac{\theta}{2\pi}\delta_Q\int_{\partial W}\chi$，这等价于要求 $\delta_Q\chi=\cF$（3d 中 $\chi$ 为 2-形式）。核对时去查 §B.4 中 B-型 3d 规范理论的 BRST 变换表，并与 p.70 的 $t=1$ 理论变换表（那里 $\delta_Q\chi=F$）对照——注意这是两个不同的理论，不能混用同一张表。第二个论证同理依赖 $\delta_Q(\eta-\ii\tilde\eta)\propto\partial_3\phi_3$；原文只说"正比于"，系数取决于约定。]

% [Tutor Note: p.62 记号不统一：拓扑项里写的是花体 $\cF$（CMSY 字体的 $\mathcal{F}$），而 p.71（$t=1$ 理论）的拓扑项与 p.70 的 BRST 表中写的是普通 $F$。按 §B.4 的约定二者都是场强，翻译时照原样保留，勿强行统一。]

沿用 [36] 中的同一思路，可以这样来描述本理论的边界条件 2-范畴：选出一个有特殊地位的边界条件 $\mathbb{X}_0$，并用任意其他边界条件 $\mathbb{X}$ 与 $\mathbb{X}_0$ 之间的缺陷线范畴 $\mathbb{V}_{\mathbb{X}\mathbb{X}_0}$ 来刻画 $\mathbb{X}$。也就是说，给每个边界条件 $\mathbb{X}$ 配上一个幺半范畴 $\mathbb{V}_{\mathbb{X}_0\mathbb{X}_0}$ 上的模范畴 $\mathbb{V}_{\mathbb{X}\mathbb{X}_0}$。

对有特殊地位的那个边界条件，一个显然的猜测是自由（Neumann）边界条件，因为它在边界上保持规范群不被破缺。要确定该边界条件下边界线算子的范畴 $\mathbb{V}_{\mathbb{X}_0\mathbb{X}_0}$，可以把 3d 理论在区间上约化，并研究所得 2d TFT 中膜的范畴。在 Neumann 情形，在区间上的约化给出如下结果：玻色场是规范场 $A$ 与 1-形式 $\phi$，费米子场是 0-形式 $\eta-\ii\tilde\eta$、1-形式 $\psi+\ii\tilde\psi$ 与 2-形式 $\chi$。这正是 2d B-型拓扑规范理论的场含量，见 §A.2。容易检验这些场的 BRST 变换也与 B-型 2d 规范理论中的相同。该 2d TFT 的膜范畴是 $G=U(1)$ 的有限维分次表示的范畴，细节见 §A.2。这是因为所能附加的边界自由度只能由携带规范群表示的向量空间来描述。幺半结构无法从 2d 考虑确定，但容易看出它由通常的张量积给出。事实上，如 §A.2 所描述，对应于表示空间 $V$ 的膜是这样得到的：把复联络 $\cA=A+\ii\phi$ 在表示 $V$ 中的和乐插入路径积分。
